% ============================================================
%  Setmana 11 — Unitat 10 · Representació de funcions
%  Exercicis: 37, 38, 43, 45, 62, 63
% ============================================================

\setmana{11}{Unitat 10 \textperiodcentered{} Representació de funcions}{29 de novembre de 2026}%
        {37, 38, 43, 45, 62 i 63}

\begin{errada}{Exercici 37}
  \ondiu{«$f'(x)>0$ en $(-\infty,-2)\cup(2,6)$» \; i \;
    «$f'(x)<0$ en $(-2,2)\cup(6,+\infty)$»}
  \hadedir{«$f'(x)>0$ en $(-\infty,x_1)\cup(x_2,x_3)$» \; i \;
    «$f'(x)<0$ en $(x_1,x_2)\cup(x_3,+\infty)$»}
  \motiu{El signe de $f'$ canvia als punts on la gràfica talla l'eix $X$, que
    l'enunciat anomena $x_1$, $x_2$ i $x_3$ i que la mateixa solució fa servir
    a la línia següent. A la gràfica, $x_1$ és entre $-2$ i $-1$ i $x_3$ entre
    $5$ i $6$: no són $-2$ ni $6$.}
\end{errada}

\begin{errada}{Exercici 38 (enunciat)}
  \ondiu{«Creixement: $(-\infty,0)\cup(2\;-\infty)$.»}
  \hadedir{«Creixement: $(-\infty,0)\cup(2,+\infty)$.»}
  \motiu{És una errada de l'enunciat. La funció creix abans del màxim $(0,0)$
    i després del mínim $(2,4)$, com es veu a la gràfica de la solució.}
\end{errada}

\begin{errada}{Exercici 45 b)}
  \ondiu{$\mathrm{Dom}\,f=\R-\{3\}$}
  \hadedir{$\mathrm{Dom}\,f=\R-\{-3\}$}
  \motiu{El denominador és $x+3$, que s'anul·la en $x=-3$.}
\end{errada}

\begin{errada}{Exercici 62 c)}
  \ondiu{$\dlim_{x\to-\infty}\,(x^4+2x^3-2x-1)=-\infty$}
  \hadedir{$\dlim_{x\to-\infty}\,(x^4+2x^3-2x-1)=+\infty$}
  \motiu{El terme dominant és $x^4$, que és positiu tant si $x\to-\infty$ com
    si $x\to+\infty$.}
\end{errada}

\begin{errada}{Exercici 62 d)}
  \ondiu{$\dlim_{x\to+\infty}\,(-2x^4+3x^2-x-1)=+\infty$}
  \hadedir{$\dlim_{x\to+\infty}\,(-2x^4+3x^2-x-1)=-\infty$}
  \motiu{El terme dominant és $-2x^4$, que tendeix a $-\infty$ tant si
    $x\to-\infty$ com si $x\to+\infty$.}
\end{errada}

\begin{errada}{Exercici 63 c)}
  \ondiu{$\dlim_{x\to-\infty}\dfrac{2x}{x+1}=2$ \; i \;
    $\dlim_{x\to+\infty}\dfrac{2x}{x+1}=2$}
  \hadedir{$\dlim_{x\to-\infty}\dfrac{4x^2}{x^2-9}=4$ \; i \;
    $\dlim_{x\to+\infty}\dfrac{4x^2}{x^2-9}=4$}
  \motiu{Els límits que hi ha impresos són d'una altra funció. Amb la de
    l'apartat, $y=\dfrac{4x^2}{x^2-9}$, el límit és $4$, que és el que dona
    l'asímptota horitzontal $y=4$ de la mateixa solució.}
\end{errada}

%\senseerrades{43}
